\section{Introduction}\label{sec:introduction}

A numerical semigroup is a cofinite additive submonoid of
\(\mathbb N=\{0,1,2,\ldots\}\). Its embedding dimension \(e\) is the
cardinality of its minimal generating set, its multiplicity \(m\) is its
least positive element, and its conductor \(c\) is the least integer such
that \(c+\mathbb N\) is contained in the semigroup. Write
\(n=|S\cap[0,c)|\). Its genus is \(g=|\mathbb N\setminus S|\),
so \(n=c-g\). Wilf's question \cite{W} asks whether
\[
en\ge c
\]
for every numerical semigroup. Fr\"oberg, Gottlieb, and H\"aggkvist
\cite{FGH} established the inequality in embedding dimension at most
three. This paper addresses the next dimension.

\begin{theorem}\label{thm:main}
Every numerical semigroup with four minimal generators satisfies
\[
W_4(S):=4|S\cap[0,c)|-c\ge0.
\]
\end{theorem}

\subsection{Foundational work}

The proof builds on Zhai's \cite{Z} preferred factorizations of
Ap\'ery elements and weighted mean inequality for finite lower ideals.
Hellus, Rechenauer, and Waldi \cite{HRW} developed the monomial
initial-ideal and residue-lattice interpretation, including restrictions
on the supports of minimal exclusions. The classical Ap\'ery genus
formula supplies the exact link from
weighted moments to Wilf's inequality; see also \cite{Book}.

Two published computer-assisted results dispose of multiplicities through
19: Bruns, Garc\'ia-S\'anchez, O'Neill, and Wilburne
\cite[Theorem~4.3]{B}, and Kliem and Stump
\cite[Proposition~6.9, published version]{KS}.
The clique-tree ingredient in our geometric decomposition is classical
\cite{G}; its needed finite form is proved here. Connections with
three-dimensional L-shapes \cite{L,C} and other approaches to Wilf's
problem are discussed in Section~\ref{sec:related}.

\subsection{General approach}

Write \(S=\langle m,a_1,a_2,a_3\rangle\), minimally. Choose one
lexicographically preferred factorization of each Ap\'ery element by
\(a_1,a_2,a_3\). Their exponent vectors form a lower ideal
\(T\subseteq\mathbb N^3\) of cardinality \(m\), with labels distinct
modulo \(m\). Set
\[
A=\min_i a_i,\qquad w_i=a_i/A,\qquad
M=c+m-1,\qquad H=M/A,\qquad s=\sum_{x\in T}x.
\]
The genus formula gives
\[
mW_4(S)=A(3mH-4w\cdot s)-m(m-1).
\]
Thus the weighted mean baseline \(3mH-4w\cdot s\ge0\) does not suffice:
the proof needs a positive surplus that pays for the correction term.

We pursue this surplus through geometry and residue collisions. A minimal
excluded vector and its representative have disjoint supports, so there
is at most one minimal exclusion with all coordinates positive. An ideal
without such an exclusion is a central anchored box with at most three
monotone rectangular horns. An ideal with one is obtained from such a
completion by deleting a single upper orthant. This decomposition makes
additive extremal bounds accessible to finite, exact dynamic programming.

For small multiplicity, an analytic reduction places every hypothetical
counterexample inside an explicit bounded generator box. For large
multiplicity, the degree of the possible positive corner and the degree
of \(T\) give an exhaustive partition. The corner \((1,1,1)\) is handled
by a counting and projection argument. Other cases use local centroid
gains, axis witnesses, or interval domination of all real weights in a
compact failure region. A finite degree strip also yields a continuous
moment gap by averaging translated lattice meshes.

The weakest retained conclusion is
\(3mH-4w\cdot s\ge m-29/10\) for the branches not already settled
directly. For \(m\ge30\), the exact moment identity and \(A\ge m+1\)
then give \(W_4>-1\); integrality gives \(W_4\ge0\).

\subsection{Organization}

Section~\ref{sec:proof} gives the proof and its seven finite assertions.
Section~\ref{sec:analytic} supplies three longer analytic arguments integral
to that proof. Section~\ref{sec:verification} specifies the computational
artifact and its checks. The final two mathematical sections discuss
extensions beyond four generators and related work. We conclude with an
account of how the proof was constructed using AI and exact computation.
